\section{Small oscillatory metric increments}\label{sec:small-modes}

We first construct a correction whose amplitude and wavelength are both small.
A good phase, in the sense of Section~\ref{sec:preliminaries}, permits an
oscillation in the normal direction perpendicular to its transverse second
form.  Finite differential corrections make this oscillation solve the
homogeneous linearized equation to high order.  Its leading quadratic mean
is a positive rank-one tensor.  Combining finitely many such oscillations
will add a prescribed positive tensor to finite accuracy while changing
the map by a quantity small in $C^2$.

We carry out the construction for the metric map and for a finite
polynomial differential perturbation of it.  In
Section~\ref{sec:primitive-realization}, removing the large velocity loop's
oscillating errors leaves mean terms that are differential operators of
the slow map, multiplied by powers of the fast scale.  Here we allow
those terms as part of the operator being corrected.
The perturbation changes the required relation between scales.  We will
bound it using a power of the short scale that is independent of the
requested finite accuracy.

\subsection{Weighted bounds and the increment statement}

For a smooth function, vector field, or tensor on the fixed coordinate atlas,
set
\[
 |f|_{m,s}=\sum_{j=0}^m s^j\|f\|_{C^j},\qquad 0<s\leq1.
\]
For a map $G$, the notation $|j^2G|_{m,s}$ means
$\sum_{j=0}^m s^j\|G\|_{C^{j+2}}$. We write $f=O_s(b)$ if, for every fixed
$m$, $|f|_{m,s}\leq C_m b$. In statements concerning a family of inputs, the
constants may depend on its indicated weighted bounds and on the fixed
coordinate data. Product estimates and smooth composition estimates hold in
these norms by the product and chain rules. In particular,
\begin{equation}\label{eq:weighted-derivative}
 |\partial^rG|_{m,s}\leq C_m s^{2-r}
 \quad(r\geq2),\qquad\text{if }j^2G=O_s(1).
\end{equation}

Fix positive phase data from Lemma~\ref{lem:phase-data}: phases $\phi_\nu$,
cutoffs $\psi_\nu$, and linear forms $q_\nu$.  Choose concentric $C^0$ balls
$\mathcal U_0,\mathcal U_1$ about a metric $h_0$, with radii
$0<r_0<r_1$, so that for $A\in\mathcal U_1$,
\begin{equation}\label{eq:small-positive-decomposition}
 b_\nu(A)=\psi_\nu\sqrt{q_\nu(A)},\qquad
 A=\sum_\nu b_\nu(A)^2\,d\phi_\nu^2.
\end{equation}
All square-root arguments have a uniform positive lower bound for
$A\in\mathcal U_1$ on their coordinate neighborhoods.  The prescribed
tensor belongs to $\mathcal U_0$; the larger ball allows the finite
intermediate tensor adjustments below.
The differentials of $\phi_\nu$, $2\phi_\nu$, and
$\phi_\nu\pm\phi_\mu$ are good wherever the corresponding supports meet.
For each such phase choose local coordinates with that phase as the first
coordinate. All these finitely many coframes and their inverses have fixed
bounds. We require the relevant immersion 2-jets to remain in a compact subset
of the open set where these conditions hold.

A differential operator is called \emph{polynomial above order two} if each of
its coordinate components is a finite sum of terms
\begin{equation}\label{eq:polynomial-jet-monomial}
 c(x,j^2G)\prod_{\ell=1}^d\partial^{\alpha_\ell}G^{k_\ell},
 \qquad |\alpha_\ell|>2,
\end{equation}
where $c$ is smooth; the empty product is allowed. This definition permits
arbitrary smooth dependence on the low jets within the prescribed open set.

\begin{proposition}[Small increment]\label{prop:small-increment}
Let the phase data above be fixed. Let $P_1,\ldots,P_d$ be fixed smooth
tensor-valued differential operators polynomial above order two, and put
\[
 \mathcal P(G)=\gamma(G)+\sum_{i=1}^d\epsilon^i P_i(G),
 \qquad 0\leq\epsilon\leq1.
\]
There is a finite exponent $p$, depending on these operators and independent of
the integer $q$ below, with the following property. Suppose $G$ is an immersion
whose 2-jets satisfy the stated compact margins, $H\in\mathcal U_0$, and
\[
 j^2G=O_s(1),\qquad H=O_s(1).
\]
For every positive integer $q$, every $\delta>0$, and scales
$0<\tau\leq s\leq1$, define
\[
 \eta=\frac\tau s+\epsilon\tau^{-p}.
\]
If $\eta$ and $\delta/\tau$ are sufficiently small, there is a smooth
displacement $U$ such that
\begin{equation}\label{eq:small-increment-estimate}
 U=O_\tau(\delta\tau),\qquad
 \mathcal P(G+U)-\mathcal P(G)-\delta^2H
   =O_\tau\left(\delta\eta^q+\frac{\delta^3}{\tau}\right).
\end{equation}
The smallness thresholds and every fixed output norm use only finitely many
weighted input bounds. Their required orders may increase with $q$.
Each local summand of $U$ is supported in an original phase support or an
intersection of two such supports. If the perturbation sum is absent, take
$\eta=\tau/s$.
\end{proposition}

The proof has three finite steps: solve the nonzero linear modes to high order,
adjust their quadratic mean, and cancel the remaining quadratic modes.
There is no infinite series of differential operators in this construction.

\subsection{The conjugated linearized equation}

Fix one permitted phase and call its coordinates $x,y$. Complexify all vector
spaces linearly. Let $g=\gamma(G)$ and $D_G=\gamma'(G)$. Decompose a complex
amplitude $Z$ as
\[
 Z=G_k g^{k\ell}u_\ell+v,
 \qquad u_i=\langle G_i,Z\rangle,
 \qquad v\perp\operatorname{span}(G_x,G_y).
\]
Since $G_{ij}=\Gamma^k_{ij}G_k+B_{ij}$, direct differentiation gives
\[
 (D_GZ)_{ij}
 =\partial_i u_j+\partial_j u_i-2\Gamma^k_{ij}u_k
       -2\langle B_{ij},v\rangle.
\]
Here and in this subsection $B=B_G$. This identity already accounts for the
normal component; no derivative of the normal frame is omitted. Conjugating by
$e^{ix/\tau}$ gives
\begin{equation}\label{eq:conjugated-metric-linearization}
 (D_G^xZ)_{ij}
 =\partial_i u_j+\partial_j u_i-2\Gamma^k_{ij}u_k
 +\frac{i}{\tau}(\delta_{ix}u_j+\delta_{jx}u_i)
 -2\langle B_{ij},v\rangle.
\end{equation}

The good-phase condition says that $b=B_{yy}$ does not vanish. Given a target
amplitude $h$ and a normal amplitude $v_0$ perpendicular to $b$, define
\begin{equation}\label{eq:normal-elimination}
 c_{ij}=\frac{\langle B_{ij},b\rangle}{|b|^2},\qquad
 v=v_0+\frac b{|b|^2}
 \left(\partial_yu_y-\Gamma^k_{yy}u_k-\frac{h_{yy}}2\right).
\end{equation}
The $yy$ equation in $D_G^xZ=h$ is then exact. The $xx$ equation, divided by 2,
and the $xy$ equation become
\begin{align}
 \frac i\tau u_x+\partial_xu_x-\Gamma^k_{xx}u_k
  -c_{xx}(\partial_yu_y-\Gamma^k_{yy}u_k)
 &=\frac{h_{xx}}2+\langle B_{xx},v_0\rangle
                      -\frac{c_{xx}h_{yy}}2,
 \label{eq:conjugated-xx}\\
 \frac i\tau u_y+\partial_xu_y+\partial_yu_x-2\Gamma^k_{xy}u_k
  -2c_{xy}(\partial_yu_y-\Gamma^k_{yy}u_k)
 &=h_{xy}+2\langle B_{xy},v_0\rangle-c_{xy}h_{yy}.
 \label{eq:conjugated-xy}
\end{align}
Thus the remaining system is
\[
 \frac i\tau u+Lu=k(h,v_0),
\]
where $L$ is first order and $k$ is algebraic and linear. All coefficients are
smooth functions of the 2-jet of $G$ on the permitted compact set.

Define $S(h,v_0)$ by taking $u=(\tau/i)k(h,v_0)$ and reconstructing $v,Z$ from
\eqref{eq:normal-elimination}. Weighted differentiation gives
\begin{align}
 |S(h,v_0)|_{m,s}
 &\leq C_m\bigl(|h|_{m+1,s}+|v_0|_{m+1,s}\bigr),
 \label{eq:first-parametrix-bound}\\
 |D_G^xS(h,v_0)-h|_{m,s}
 &\leq C_m\frac\tau s
      \bigl(|h|_{m+1,s}+|v_0|_{m+1,s}\bigr).
 \label{eq:first-parametrix-error}
\end{align}
Indeed the residual is $Lu$, and the reconstructed normal component contains
at most one derivative of $u$. In the homogeneous case the leading term is
exactly the chosen normal amplitude:
\begin{equation}\label{eq:free-mode-leading}
 |S(0,v_0)-v_0|_{m,s}
 \leq C_m\frac\tau s|v_0|_{m+1,s}.
\end{equation}
These are differential formulas, so they preserve supports.

\subsection{A fixed exponent for polynomial perturbations}

We next quantify the effect of the operators $P_i$. This is the point at which
independence from the approximation order $q$ is needed.

For a monomial \eqref{eq:polynomial-jet-monomial}, let
\[
 E_\alpha=\sum_{\ell=1}^d(|\alpha_\ell|-2),
 \qquad E=\max E_\alpha,
\]
where the maximum runs over all monomials of all $P_i$. For an empty collection
take $E=0$. A permissible choice will be
\begin{equation}\label{eq:fixed-perturbation-power}
 p=E+6.
\end{equation}

To verify this, consider a variation of order $k\leq3$. Suppose $h$ arguments
replace higher-jet factors of orders $r_1,\ldots,r_h$ and the other $k-h$
arguments differentiate the smooth low-jet coefficient. The unvaried factors
have excess weight $E_\alpha-\sum_{j=1}^h(r_j-2)$. The variation arguments
carry at most $\sum_jr_j+2(k-h)$ derivatives. The total short-scale power is
therefore at most
\[
 E_\alpha-\sum_{j=1}^h(r_j-2)+\sum_{j=1}^hr_j+2(k-h)
 =E_\alpha+2k\leq E+6.
\]
Weighted differentiation does not increase this exponent:
\eqref{eq:weighted-derivative} controls each differentiated higher-jet factor,
and the chain rule controls the smooth low-jet coefficient. For conjugated
amplitudes, each derivative of an exponential costs a power of $\tau^{-1}$,
while a slow derivative costs $s^{-1}\leq\tau^{-1}$.

Put $\mathcal R=\mathcal P-\gamma$. For $1\leq k\leq3$, this proves, with a
fixed finite derivative loss $\ell$, the bounds
\begin{align}
 &\left|e^{-i(\phi_1+\cdots+\phi_k)/\tau}
 \mathcal R^{(k)}(G)
 \bigl(e^{i\phi_1/\tau}Z_1,\ldots,e^{i\phi_k/\tau}Z_k\bigr)
 \right|_{m,s}
 \notag\\[-1mm]
 &\hspace{35mm}\leq
 C_m\epsilon\tau^{-p}\prod_{j=1}^k|Z_j|_{m+\ell,s}.
 \label{eq:conjugated-perturbation-bound}
\end{align}
The phases in this estimate belong to the fixed finite collection. At scale
$\tau$, the same estimate without exponentials holds for arbitrary variations
at every map $K$ with $j^2K=O_\tau(1)$ in the permitted low-jet domain.
The constants use finitely many corresponding weighted bounds. The inequality
$\epsilon^i\leq\epsilon$ handles all the perturbation coefficients.

Let $\mathcal L=(\mathcal P'(G))^x$ and $S_0(h)=S(h,0)$. Combining
\eqref{eq:first-parametrix-bound}--\eqref{eq:first-parametrix-error} with
\eqref{eq:conjugated-perturbation-bound} gives
\[
 |\mathcal L S(h,v_0)-h|_{m,s}
 \leq C_m\eta\bigl(|h|_{m+\ell,s}+|v_0|_{m+\ell,s}\bigr),
\]
after increasing $\ell$. Starting with $Z_0=S(h,v_0)$, define
\[
 r_j=\mathcal LZ_j-h,\qquad Z_{j+1}=Z_j-S_0(r_j).
\]
The exact identity
\[
 r_{j+1}=-(\mathcal L S_0-\Id)r_j
\]
shows by finite induction that, for any integer $q$, the resulting amplitude
$Z$ satisfies
\begin{align}
 |Z|_{m,s}&\leq C_{m,q}
       \bigl(|h|_{m+\ell_q,s}+|v_0|_{m+\ell_q,s}\bigr),
 \label{eq:finite-mode-bound}\\
 |\mathcal LZ-h|_{m,s}&\leq C_{m,q}\eta^q
       \bigl(|h|_{m+\ell_q,s}+|v_0|_{m+\ell_q,s}\bigr).
 \label{eq:finite-mode-error}
\end{align}
For $h=0$ it also satisfies
\begin{equation}\label{eq:finite-free-mode}
 |Z-v_0|_{m,s}\leq C_{m,q}\eta|v_0|_{m+\ell_q,s}.
\end{equation}
The integer $\ell_q$ and the constants may increase with $q$, but the exponent
$p$ in \eqref{eq:fixed-perturbation-power} has already been fixed. Every step is
differential and support preserving.

We have now solved every allowed nonzero mode to any fixed power of $\eta$.
It remains to select the free amplitudes so that their quadratic mean is the
desired tensor and then remove the oscillatory quadratic error.

\subsection{The quadratic mean and its finite adjustment}

Finite inner corrections to oscillatory amplitudes also occur in
K\"all\'en's low-regularity embedding construction
\cite[Lemma~3.3 and its proof]{Kallen1978}.
We derive the mean equation and its finite substitution estimates for
the differential operators considered here.

In a phase chart choose a smooth unit normal $n_\nu$ perpendicular to
$(B_G)_{yy}$. One explicit choice is the normalized ambient perpendicular
product of $G_x,G_y,(B_G)_{yy}$. The vectors in this product are independent.
The choice depends smoothly on the 2-jet and uses only the orientation of this
chart and of $\R^4$.

For a trial tensor $A\in\mathcal U_1$, use
\[
 v_{0,\nu}=\sqrt2\,\delta\tau b_\nu(A)n_\nu,\qquad h=0,
\]
in the finite mode construction, and denote its amplitude by $Z_\nu(A)$.
Define the real displacement
\[
 U_1(A)=\sum_\nu\operatorname{Re}
             \bigl(e^{i\phi_\nu/\tau}Z_\nu(A)\bigr).
\]
Then $U_1=O_\tau(\delta\tau)$. Before multiplication by the exponential, the
leading complex amplitude of $dU_1/\delta$ in the $\nu$th mode is
\[
 i\sqrt2\,b_\nu(A)n_\nu\,d\phi_\nu,
\]
with error $O_s(\eta)$. This follows from \eqref{eq:finite-free-mode}; a
slow derivative of the leading amplitude contributes $O_s(\tau/s)$.
Furthermore,
\begin{equation}\label{eq:free-sum-linear-error}
 \mathcal P'(G)U_1=O_\tau(\delta\tau\eta^q).
\end{equation}

Expand the quadratic tensor
\[
 Q(A)=\tfrac12\mathcal P''(G)(U_1(A),U_1(A))
\]
using the positive and negative mode labels. Its \emph{zero-phase part}, denoted
$Q_0(A)$, consists of a mode paired with its own complex conjugate. All other
terms have phases $\pm2\phi_\nu$ or $\pm\phi_\nu\pm\phi_\mu$ with
$\nu\ne\mu$. The latter phases have nonzero differentials on their supports,
by hypothesis; in particular they do not contribute another zero-phase term.
Repeated nonzero phases may simply be grouped.

For one leading real mode, the metric quadratic term is
\[
 2\delta^2b_\nu(A)^2\sin^2(\phi_\nu/\tau)\,d\phi_\nu^2
 =\delta^2b_\nu(A)^2
       \bigl(1-\cos(2\phi_\nu/\tau)\bigr)\,d\phi_\nu^2.
\]
Its zero-phase term is $\delta^2b_\nu(A)^2d\phi_\nu^2$. Thus
\eqref{eq:small-positive-decomposition} and the amplitude error give
\begin{equation}\label{eq:quadratic-mean-map}
 \delta^{-2}Q_0(A)=A+T(A),\qquad T(A)=O_s(\eta).
\end{equation}
The contribution of $\mathcal R''$ has the same bound, since
\eqref{eq:conjugated-perturbation-bound} bounds it by
$C\epsilon\tau^{-p}(\delta\tau)^2\leq C\delta^2\eta$.
Every nonzero-mode coefficient of $Q$ is $O_s(\delta^2)$.

The map $T$ has a finite derivative-loss difference estimate:
\begin{equation}\label{eq:mean-map-difference}
 |T(A)-T(\bar A)|_{m,s}
 \leq C_m\eta|A-\bar A|_{m+L_q,s}
\end{equation}
on bounded sets of weighted norms in $\mathcal U_1$. To see
this, the square roots in $b_\nu$ have positive lower bounds, so their
chain-rule difference estimates apply. The finite mode construction is linear
in $b_\nu(A)n_\nu$. In the difference of the two sides of
\eqref{eq:quadratic-mean-map}, the leading contribution $A-\bar A$ cancels
exactly. Every remaining bilinear term contains either one $O_s(\eta)$ mode
error or the factor $\epsilon\tau^{-p}$ from the perturbation. The product
rule then gives \eqref{eq:mean-map-difference} with a finite $L_q$.

Set $A_0=H$ and successively define
\[
 A_{j+1}=H-T(A_j).
\]
For any fixed number of substitutions, all these tensors remain in
$\mathcal U_1$ and have bounded weighted norms when $\eta$ is sufficiently
small.  The positive gap $r_1-r_0$ and the estimate $T(A_j)=O_s(\eta)$
give this domain condition using only finitely many input derivatives:
for a fixed number of substitutions, work backward through their finite
derivative losses to obtain the required $C^0$ bounds.  Once all iterates
lie in $\mathcal U_1$, the finite formulas give bounds in each higher
weighted norm by the chain rule, without a further domain restriction.
Their successive
differences gain one power of $\eta$ by
\eqref{eq:mean-map-difference}. More explicitly, for $j\geq1$,
\[
 A_j+T(A_j)-H=T(A_j)-T(A_{j-1})=O_s(\eta^{j+1}).
\]
After finitely many substitutions, choose the resulting $A$ so that
\begin{equation}\label{eq:adjusted-quadratic-mean}
 Q_0(A)=\delta^2H+O_s(\delta^2\eta^q).
\end{equation}
This is a finite algebraic procedure on smooth tensors, not a contraction
argument in a function space with derivative loss.

\subsection{Cancellation and the cubic remainder}

Use the amplitude $A$ just obtained. On each nonzero quadratic phase, apply
\eqref{eq:finite-mode-bound}--\eqref{eq:finite-mode-error} with target the
negative corresponding coefficient of $Q-Q_0$ and with $v_0=0$. Taking real
parts gives $U_2$ such that
\begin{equation}\label{eq:quadratic-cancellation}
 U_2=O_\tau(\delta^2),\qquad
 \mathcal P'(G)U_2+(Q-Q_0)=O_\tau(\delta^2\eta^q).
\end{equation}
The supports are intersections of the original supports, or original supports
for doubled phases. Since $\delta/\tau$ is small,
\[
 U=U_1+U_2=O_\tau(\delta\tau),\qquad
 dU_1=O_\tau(\delta),\qquad dU_2=O_\tau(\delta^2/\tau).
\]
The metric quadratic terms involving $U_2$ are bounded by
\[
 C_m\left(\frac{\delta^3}{\tau}+\frac{\delta^4}{\tau^2}\right)
 \leq C_m\frac{\delta^3}{\tau}.
\]
The corresponding terms from $\mathcal R''$ are, by
\eqref{eq:conjugated-perturbation-bound} at scale $\tau$, bounded by
\[
 C_m\epsilon\tau^{-p}(\delta^3\tau+\delta^4)
 \leq C_m\frac{\delta^3}{\tau}.
\]
Here $\epsilon\tau^{-p}\leq\eta\leq1$ and $\delta/\tau\leq1$ suffice.

For completeness, the third-order Taylor estimate is uniform along the whole
segment $G+tU$, $0\leq t\leq1$. Indeed
\[
 |j^2U|_{m,\tau}\leq C_m\left(\frac \delta\tau+\frac{\delta^2}{\tau^2}\right),
\]
so smallness of $\delta/\tau$ keeps the low jets inside the permitted domain.
Also $j^2(G+tU)=O_\tau(1)$, since $\tau\leq s$. The arbitrary-variation
version of \eqref{eq:conjugated-perturbation-bound} therefore bounds the integral
Taylor remainder of $\mathcal R$ by
\[
 C_m\epsilon\tau^{-p}(\delta\tau)^3
 \leq C_m\frac{\delta^3}{\tau}.
\]
There is no third-order remainder for $\gamma$ itself.

Finally combine \eqref{eq:free-sum-linear-error},
\eqref{eq:adjusted-quadratic-mean}, and \eqref{eq:quadratic-cancellation}.
Their errors are bounded by $C_m \delta\eta^q$, because $\delta,\tau\leq1$.
The remaining terms have just been bounded by $C_m\delta^3/\tau$. This proves
\eqref{eq:small-increment-estimate} and Proposition~\ref{prop:small-increment}.
Every local formula is smooth and differential in cutoff amplitudes, so
extension by zero gives the asserted global smooth displacement and support
properties. None of the local normal choices requires orientability of $M$.

\begin{remark}[Dependence on geometric bounds]\label{rem:small-polynomial-control}
For $\mathcal P=\gamma$, the constants and inverse smallness thresholds at any
fixed accuracy and fixed output order can be bounded polynomially in finitely
many input and phase-data bounds and reciprocal nondegeneracy margins.
Indeed the coordinate construction uses metric inversion, contractions,
division by $|B_{yy}|^2$, and normalized perpendicular products. Derivatives of
these expressions have polynomial bounds in their entries, derivatives, and
reciprocal denominators. The square roots in
\eqref{eq:small-positive-decomposition} have the same property on a positive
input interval. The finite residual corrections and mean substitutions use
only a fixed number of these operations. Summing over a varying finite phase
collection adds a polynomial dependence on its cardinality. This observation
will allow the geometric bounds to depend polynomially on the parameter of an
approximate immersion family.
\end{remark}
